\documentclass[11pt,a4paper]{amsart}
\usepackage[T1]{fontenc}
\usepackage{lmodern}
\usepackage[margin=25mm]{geometry}
\usepackage{amsmath,amssymb,amsthm}
\usepackage[expansion=false]{microtype}
\usepackage{array,booktabs,tabularx}
\usepackage{needspace}
\setlength{\captionindent}{0pt}
\usepackage{xurl}
\usepackage[hidelinks]{hyperref}
\urlstyle{same}
\hypersetup{
  pdftitle={Erd\H{o}s Problem 928: quantitative comparisons and original formulations},
  pdfauthor={JI HO BAE},
  pdfsubject={Elementary comparison estimates and the deduction of Erd\H{o}s Problem 928 from OpenAI's joint Dickman law}
}
\newtheorem{proposition}{Proposition}
\newtheorem{corollary}[proposition]{Corollary}
\newcommand{\N}{\mathbb N}
\newcommand{\A}{\mathcal A}
\newcommand{\B}{\mathcal B}
\newcommand{\C}{\mathcal C}
\newcolumntype{L}[1]{>{\raggedright\arraybackslash}p{#1}}
\newcolumntype{Y}{>{\raggedright\arraybackslash}X}
\title[Erd\H{o}s Problem 928 and the joint Dickman law]{Erd\H{o}s Problem 928:\\
quantitative comparisons and original formulations}
\author{JI HO BAE}
\date{October 9, 2026}

\begin{document}
\begin{abstract}
We show that OpenAI's joint Dickman law yields the density
$\rho(1/\alpha)\rho(1/\beta)$ under the conditions stated in
Erd\H{o}s Problem~928. Here $\rho$ is the Dickman--de Bruijn function.
For fixed $0<\alpha,\beta<1$ and real $X\ge2$, we compare the sets of
integers specified by the problem and the theorem in the range
$2\le n\le X$. Their symmetric difference has size at most
\mbox{$\pi(X^\alpha)+\pi((X+1)^\beta)$}, where $\pi$ counts primes.
This elementary bound handles both the strict inequalities and the
shift of the second threshold from $n^\beta$ to $(n+1)^\beta$.
We also establish the density for Erd\H{o}s's strict formulation in
which both thresholds are powers of $n$, and compare the statements
in Erd\H{o}s's two 1976 papers and the 1978 Erd\H{o}s--Pomerance paper.
The analytic joint law is due to OpenAI. This note gives explicit
finite comparison bounds and shows how the original formulations are related.
\end{abstract}
\maketitle

\section{Introduction}

For an integer $n\ge2$, let $P(n)$ be the largest prime factor of $n$, and set
$P(1)=1$. Write $\rho$ for the Dickman--de Bruijn function: the continuous
function on $[0,\infty)$ determined by
\[
 \rho(u)=1\quad(0\le u\le1),\qquad
 u\rho'(u)=-\rho(u-1)\quad(u>1).
\]
The classical smooth-number limit is
\begin{equation}\label{eq:marginal}
 \lim_{X\to\infty}\frac1X
 \#\{1\le n\le X:P(n)\le X^c\}=\rho(1/c)
 \qquad(0<c<1);
\end{equation}
see de Bruijn~\cite[p.~51, (1.1)--(1.2)]{deBruijn1951}.
We also use the properties $\rho(u)>0$ for finite $u\ge0$
and $\rho(u)\to0$ as $u\to\infty$; see
de Bruijn~\cite[p.~52, (1.8)]{deBruijn1951} and the defining equation.

Throughout, $\N=\{1,2,\ldots\}$. Counting variables are integers,
and the cutoff $X$ is real. Natural density means the limit of the
unweighted count up to $X$, divided by $X$.

Problem~928 in Bloom's collection~\cite{Bloom928} asks whether the
integers satisfying
\begin{equation}\label{eq:problem}
 P(n)<n^\alpha,\qquad P(n+1)<(n+1)^\beta
\end{equation}
have a natural density for every fixed $\alpha,\beta\in(0,1)$.
The accompanying discussion asks whether that density equals
$\rho(1/\alpha)\rho(1/\beta)$. The cited sources are Erd\H{o}s's
Manitoba and Debrecen papers~\cite{Erdos1976Manitoba,Erdos1976Debrecen}
and the paper of Erd\H{o}s and Pomerance~\cite{ErdosPomerance1978}.
We compare their statements in Section~\ref{sec:sources}.

OpenAI's joint Dickman law~\cite[Theorem~1.1, p.~2]{OpenAI2026}
states that, for every fixed
$\alpha,\beta\in(0,1)$,
\begin{equation}\label{eq:JD}
 \lim_{X\to\infty}\frac1X
 \#\{2\le n\le X:P(n)\le n^\alpha,\ P(n+1)\le n^\beta\}
 =\rho(1/\alpha)\rho(1/\beta).
\end{equation}
OpenAI states that this theorem resolves the Erd\H{o}s--Pomerance
joint Dickman conjecture. Section~11.3 of the same paper also gives
the versions with thresholds $X^\alpha,X^\beta$ and the corresponding
upper tails.

To pass from~\eqref{eq:JD} to~\eqref{eq:problem}, one must make the
inequalities strict and base the second threshold on the second
integer of the pair. Proposition~\ref{prop:comparison} bounds the
resulting discrepancy and also treats Erd\H{o}s's strict formulation.
Section~\ref{sec:sources} identifies the original statements and
explains how their inequalities and counting ranges are related.

\section{Finite comparison estimates}\label{sec:comparison}

Fix $\alpha,\beta\in(0,1)$. For $X\ge2$, define
\begin{align*}
 \A_{\alpha,\beta}(X)
   &=\{n\in\N:2\le n\le X,\ P(n)<n^\alpha,
                                \ P(n+1)<(n+1)^\beta\},\\
 \B_{\alpha,\beta}(X)
   &=\{n\in\N:2\le n\le X,\ P(n)\le n^\alpha,
                                \ P(n+1)\le n^\beta\},\\
 \C_{\alpha,\beta}(X)
   &=\{n\in\N:2\le n\le X,\ P(n)<n^\alpha,
                                \ P(n+1)<n^\beta\}.
\end{align*}
Here $\A$ is the set in Problem~928, $\B$ is the set counted
in~\eqref{eq:JD}, and $\C$ is the set in condition~(1) of
Erd\H{o}s's Manitoba paper~\cite[p.~25]{Erdos1976Manitoba}.
Let
$\pi(y)=\#\{p\le y:p\text{ is prime}\}$, and write
$S\mathbin\triangle T$ for the symmetric difference of two sets.

\begin{proposition}[Comparison of conventions]\label{prop:comparison}
For every $\alpha,\beta\in(0,1)$ and every real $X\ge2$,
\begin{align}
 \#\bigl(\A_{\alpha,\beta}(X)\mathbin\triangle
          \B_{\alpha,\beta}(X)\bigr)
 &\le \pi(X^\alpha)+\pi((X+1)^\beta),\label{eq:AB}\\
 0\le\#\B_{\alpha,\beta}(X)-\#\C_{\alpha,\beta}(X)
 &\le \pi(X^\alpha)+\pi(X^\beta).\label{eq:BC}
\end{align}
\end{proposition}

\begin{proof}
We suppress the subscripts and the argument $X$ throughout the proof.
First let $n\in\B\setminus\A$. The inequality
$P(n+1)\le n^\beta<(n+1)^\beta$ shows that the second condition
for $\A$ holds, so its first condition must fail: $P(n)=n^\alpha$.
Each prime $p=P(n)\le X^\alpha$ determines at most one such integer,
namely $n=p^{1/\alpha}$. Hence
$\#(\B\setminus\A)\le\pi(X^\alpha)$.

Next let $n\in\A\setminus\B$. The first condition for $\B$ is
satisfied, so the second must fail:
\[
 n^\beta<P(n+1)<(n+1)^\beta.
\]
The intervals $(m^\beta,(m+1)^\beta)$ for integers $m\ge2$ are
pairwise disjoint, since $\beta>0$. Thus the map $n\mapsto P(n+1)$
is injective on $\A\setminus\B$. Its image consists of primes below
$(X+1)^\beta$, giving $\#(\A\setminus\B)\le\pi((X+1)^\beta)$.
Together, the two bounds prove~\eqref{eq:AB}.

Finally, $\C\subseteq\B$, and each $n\in\B\setminus\C$ satisfies
$P(n)=n^\alpha$ or $P(n+1)=n^\beta$. As above, the first equality
holds for at most $\pi(X^\alpha)$ integers. For the second,
$n=P(n+1)^{1/\beta}$ is determined by a prime at most $X^\beta$,
so at most $\pi(X^\beta)$ integers occur. Adding these bounds
proves~\eqref{eq:BC}.
\end{proof}

\begin{corollary}[Problem~928 and Erd\H{o}s's strict formulation]\label{cor:density}
For every fixed $\alpha,\beta\in(0,1)$,
\begin{equation}\label{eq:conclusion}
 \lim_{X\to\infty}\frac{\#\A_{\alpha,\beta}(X)}X
 =\lim_{X\to\infty}\frac{\#\C_{\alpha,\beta}(X)}X
 =\rho(1/\alpha)\rho(1/\beta).
\end{equation}
Both the density in Problem~928 and the density $d(\alpha,\beta)$
in Erd\H{o}s's Manitoba paper~\cite[p.~25]{Erdos1976Manitoba} therefore
exist and equal the predicted product.
\end{corollary}

\begin{proof}
For finite sets, $|\#S-\#T|\le\#(S\mathbin\triangle T)$.
By Proposition~\ref{prop:comparison} and the bound
$\pi(y)\le\lfloor y\rfloor$, both discrepancies are
$O(X^\alpha+X^\beta)=o(X)$. They therefore vanish after division by
$X$, and~\eqref{eq:JD} gives the two limits in~\eqref{eq:conclusion}.
\end{proof}

The comparison above uses no part of the joint Dickman law.
For fixed $\alpha,\beta$, it measures the error introduced by changing
the inequalities and thresholds. The limiting value is supplied
by~\eqref{eq:JD}; these estimates do not give a convergence rate for
that law.

\section{Comparison with the original sources}\label{sec:sources}

We call a threshold \emph{moving} if it depends on the integer being
counted, and \emph{fixed} if it remains constant within the counting
range. Thus $n^\alpha$ and $(n+1)^\beta$ are moving thresholds, whereas
$X^\alpha$ and $X^\beta$ are fixed thresholds. Both vary with the
counting endpoint. Put
\[
 a(0)=1,\qquad a(t)=1-\rho(1/t)\quad(0<t\le1).
\]
Table~\ref{tab:forms} compares OpenAI's theorem with the four source
formulations. Its final column gives the normalized limits and the
results that establish them. The following subsections explain the
inequalities and counting ranges in each source.

\begin{table}[htbp]
\small
\caption{Reference theorem and source formulations.
In each row, $N$ counts the integers satisfying the listed conditions.
Exponents are fixed: \mbox{$0<\alpha,\beta<1$},
\mbox{$\gamma>1$}, and \mbox{$0\le t,s\le1$}.
Here $c_\gamma=\rho(\gamma)$.
Unless indicated otherwise, \mbox{$2\le n\le X$} and $X\to\infty$.}
\label{tab:forms}
\begin{tabularx}{\textwidth}{@{}L{43mm}L{45mm}Y@{}}
\toprule
Formulation and source & Conditions & Limit and justification\\
\midrule
\textbf{Reference theorem}\newline
OpenAI~\cite[Thm.~1.1, p.~2]{OpenAI2026}
 & $P(n)\le n^\alpha$\newline
   $P(n+1)\le n^\beta$
 & $\displaystyle\frac{N}{X}\to\rho(1/\alpha)\rho(1/\beta)$\newline
   Joint Dickman law~\eqref{eq:JD}\\[5pt]
\textbf{Problem~928}\newline
Bloom~\cite{Bloom928}; Erd\H{o}s's independence statement
\cite[p.~25]{Erdos1976Manitoba}
 & $P(n)<n^\alpha$\newline
   $P(n+1)<(n+1)^\beta$
 & $\displaystyle\frac{N}{X}\to\rho(1/\alpha)\rho(1/\beta)$\newline
   This note: Proposition~\ref{prop:comparison} and
   Corollary~\ref{cor:density} (set $\A$)\\[5pt]
\textbf{Erd\H{o}s's strict formulation}\newline
Manitoba paper, condition~(1)
\cite[p.~25]{Erdos1976Manitoba}
 & $P(n)<n^\alpha$\newline
   $P(n+1)<n^\beta$
 & $\displaystyle\frac{N}{X}\to\rho(1/\alpha)\rho(1/\beta)$\newline
   This note: Proposition~\ref{prop:comparison} and
   Corollary~\ref{cor:density} (set $\C$)\\[5pt]
\textbf{Fixed-threshold pair}\newline
Erd\H{o}s (Debrecen), $r=2$
\cite[pp.~272--273]{Erdos1976Debrecen}
 & $P(n+1)\le k$\newline
   $P(n+2)\le k$\newline
   $1\le n<k^\gamma$, $k\to\infty$
 & $\displaystyle\frac{N}{k^\gamma}\to\rho(\gamma)^2=c_\gamma^2$\newline
   OpenAI~\cite[(11.7), p.~80]{OpenAI2026}, with the
   change of variables in Section~\ref{subsec:debrecen}\\[5pt]
\textbf{Fixed upper tails}\newline
Erd\H{o}s--Pomerance~\cite[p.~311]{ErdosPomerance1978}
 & $P(n)\ge X^t$\newline
   $P(n+1)\ge X^s$
 & $\displaystyle\frac{N}{X}\to a(t)a(s)$\newline
   OpenAI~\cite[(11.8) and following paragraph, p.~80]{OpenAI2026},
   including equality and endpoints\\
\bottomrule
\end{tabularx}
\end{table}

\subsection{Erd\H{o}s's Manitoba paper}
On page~25 of his Manitoba paper~\cite{Erdos1976Manitoba},
Erd\H{o}s expects the conditions $P(n)<n^\alpha$ and
$P(n+1)<(n+1)^\beta$ to be independent. His next sentence asks whether
the integers satisfying condition~(1), namely $P(n)<n^\alpha$ and
$P(n+1)<n^\beta$, have a density $d(\alpha,\beta)$.
The inequalities on Bloom's webpage are those in the independence
statement; condition~(1) instead uses $n^\beta$ as its second threshold.
Corollary~\ref{cor:density} gives the predicted density in both cases.

Positive density also answers Erd\H{o}s's question about infinitely
many solutions for each fixed pair $\alpha,\beta>0$.
For parameters in $(0,1)$, this follows from~\eqref{eq:conclusion};
for any other positive pair, choose smaller positive parameters inside
that interval. As Bloom's webpage~\cite{Bloom928} notes, infinitude
already follows from
Schinzel~\cite[Theorem~14, p.~231]{Schinzel1967}, applied to
$(2n+1)^2-1=4n(n+1)$. The theorem gives a constant $K>0$ and infinitely
many integers $n\ge3$ such that
\[
 P(n(n+1))\le n^{K/\log\log n}.
\]
Since $K/\log\log n\to0$, these integers satisfy Erd\H{o}s's strict
inequalities for every prescribed positive pair $\alpha,\beta$, once
$n$ is sufficiently large.

\subsection{Erd\H{o}s's Debrecen paper}\label{subsec:debrecen}
Write $\gamma>1$ for the exponent that Erd\H{o}s denotes by $\alpha$
in his Debrecen paper~\cite[pp.~272--273]{Erdos1976Debrecen}.
He recalls that the number of $m\le k^\gamma$ satisfying $P(m)\le k$
is $(c_\gamma+o(1))k^\gamma$, and conjectures the analogous product
formula for each fixed integer $r\ge1$: as $k$ tends to infinity
through positive integers,
\begin{equation}\label{eq:Debrecen}
 \#\{1\le n<k^\gamma:P((n+1)\cdots(n+r))\le k\}
  =(c_\gamma^r+o(1))k^\gamma.
\end{equation}
The marginal law~\eqref{eq:marginal} gives $c_\gamma=\rho(\gamma)$.

The case of two consecutive integers follows from OpenAI's
fixed-threshold law~\cite[(11.7), p.~80]{OpenAI2026}: for fixed
$c,d\in(0,1)$,
\begin{equation}\label{eq:fixed}
 \lim_{X\to\infty}\frac1X
 \#\{2\le m\le X:P(m)\le X^c,\ P(m+1)\le X^d\}
 =\rho(1/c)\rho(1/d).
\end{equation}
To apply this law to~\eqref{eq:Debrecen} with $r=2$, take $X=k^\gamma$,
$m=n+1$, and $c=d=1/\gamma$.
Since $P(m(m+1))=\max\{P(m),P(m+1)\}$, the prime-factor conditions
agree with those in~\eqref{eq:fixed}.

After the shift, the ranges $1\le n<X$ and $2\le m\le X$ differ by
at most one endpoint term. This difference disappears after
normalization, so~\eqref{eq:fixed} proves~\eqref{eq:Debrecen} for $r=2$.

Bloom's webpage~\cite{Bloom928} also gives a higher-order conjecture with
$P(n+i)\le n^\alpha$ for $0\le i<r$. Its case $r=2$ is precisely
\eqref{eq:JD} with $\beta=\alpha$. The two-variable results used here
do not prove the cases $r\ge3$.

\Needspace{10\baselineskip}
\subsection{Erd\H{o}s and Pomerance}
Erd\H{o}s and Pomerance~\cite[p.~311]{ErdosPomerance1978} express their
independence conjecture in terms of upper tails at fixed thresholds.
For $0\le t,s\le1$,
they ask whether
\begin{equation}\label{eq:upper}
 b(t,s)=\lim_{X\to\infty}\frac1X
 \#\{2\le n\le X:P(n)\ge X^t,\ P(n+1)\ge X^s\}
\end{equation}
exists, is continuous on $[0,1]^2$, and equals $a(t)a(s)$.
Their $a(t)$ denotes the corresponding one-variable upper-tail limit.
The marginal law~\eqref{eq:marginal} identifies their $a(t)$ with our
definition once we include equality and check the endpoints, as below.

For $0<t,s<1$, inclusion--exclusion applied to~\eqref{eq:fixed} and the
two marginal laws gives the limit $a(t)a(s)$ for the strict upper-tail
conditions $P(n)>X^t$, $P(n+1)>X^s$.
Replacing $n$ by $n+1$ in a marginal count changes its counting range
by only $O(1)$ terms, so its limit is unchanged.

Allowing equality leaves the limit unchanged. For fixed $0<t\le1$,
equality $P(m)=X^t$ requires $X^t$ to be a prime $p$.
Since $p\mid m$, the number of such integers $m\le X+1$ is at most
$\lfloor(X+1)/p\rfloor=O(X^{1-t}+1)=o(X)$.

At exponent $0$, the upper-tail condition holds automatically for $m\ge2$.
If just one exponent is zero, the marginal law gives the limit for the
remaining condition; if both are zero, both conditions are automatic.
At exponent $1$, the condition $P(m)\ge X$ in
the range $2\le m\le X+1$ forces $m\in[X,X+1]$, so at most two
integers occur. These observations prove~\eqref{eq:upper} on the boundary as well.
The continuity of $\rho$, together with $\rho(1)=1$ and
$\rho(u)\to0$ as $u\to\infty$, shows that $a$ is continuous on $[0,1]$.
Hence $b(t,s)=a(t)a(s)$ is continuous on the whole square.

OpenAI already gives this argument, including equality and the
endpoint cases~\cite[(11.8) and following paragraph, p.~80]{OpenAI2026}.
Here we explain how it yields the limit for the weak upper-tail
conditions used by Erd\H{o}s and Pomerance~\cite{ErdosPomerance1978}.
Together, the cited joint Dickman law, its fixed-threshold version,
and the comparisons above give all four limits in Table~\ref{tab:forms}.

\begin{thebibliography}{9}
\bibitem{Bloom928}
T.~F. Bloom, \emph{Erd\H{o}s Problem \#928},
accessed October~9, 2026.
\newline\url{https://www.erdosproblems.com/928}.

\bibitem{deBruijn1951}
N.~G. de Bruijn, \emph{On the number of positive integers $\le x$ and
free of prime factors $>y$}, Nederl. Akad. Wetensch. Proc. Ser. A
\textbf{54} (1951), 50--60.
\newline\url{https://pure.tue.nl/ws/files/1793853/597499.pdf}.

\bibitem{Erdos1976Manitoba}
P.~Erd\H{o}s, \emph{Problems and results on number theoretic properties of
consecutive integers and related questions}, Proceedings of the Fifth
Manitoba Conference on Numerical Mathematics (Winnipeg, 1975), Congressus
Numerantium XVI, Utilitas Math., Winnipeg, 1976, pp.~25--44.
\newline\url{https://users.renyi.hu/~p_erdos/1976-39.pdf}.

\bibitem{Erdos1976Debrecen}
P.~Erd\H{o}s, \emph{Problems and results on consecutive integers}, Publ.
Math. Debrecen \textbf{23} (1976), 271--282.
\newline\url{https://doi.org/10.5486/PMD.1976.23.3-4.15}.

\bibitem{ErdosPomerance1978}
P.~Erd\H{o}s and C.~Pomerance, \emph{On the largest prime factors of $n$
and $n+1$}, Aequationes Math. \textbf{17} (1978), 311--321.
\newline\url{https://users.renyi.hu/~p_erdos/1978-29.pdf}.

\bibitem{OpenAI2026}
OpenAI, \emph{The joint Dickman law for consecutive integers}, OpenAI
Math Release preprint, September~24, 2026, commit
\texttt{fd4aeeb2ee4f}.
\newline\url{https://github.com/openai/math/blob/fd4aeeb2ee4fc729c18d98444fed42fd0529eeeb/preprints/The-joint-Dickman-law-for-consecutive-integers-September-24-2026/paper.pdf}.

\bibitem{Schinzel1967}
A.~Schinzel, \emph{On two theorems of Gelfond and some of their
applications}, Acta Arith. \textbf{13} (1967), 177--236.
\newline\url{https://doi.org/10.4064/aa-13-2-177-236}.
\end{thebibliography}

\end{document}
